\documentclass[11pt,a4paper]{article}
\usepackage[utf8]{inputenc}
\usepackage[T1]{fontenc}
\usepackage{amsmath,amssymb}
\usepackage{booktabs}
\usepackage{graphicx}
\usepackage{multirow}
\usepackage{array}
\usepackage[margin=2.5cm]{geometry}
\usepackage{hyperref}
\usepackage{xcolor}
\usepackage{natbib}
\usepackage{tabularx}
\usepackage{adjustbox}

\title{CaCLR: Zero-Label Cross-Lingual Calibration via\\Correctness-Direction Alignment in Large Language Models}

\author{Anonymous}

\date{}

\begin{document}
\maketitle

%% ============================================================
%% ABSTRACT
%% ============================================================
\begin{abstract}
Non-English languages suffer substantially worse calibration in multilingual large language models, yet all existing probe-based calibration methods operate exclusively in English.
We present the first systematic study of cross-lingual correctness-probe transfer across 20 languages on two large-scale architectures (Qwen3-235B-A22B and Llama~4 Scout 109B).
We discover an architecture-dependent \emph{depth--transferability Pareto frontier}: English-optimal probe layers consistently degrade in cross-lingual transfer on Qwen3 ($\rho \approx -0.65$) but show a weaker, language-heterogeneous tradeoff on Llama~4 ($\rho \approx -0.33$).
We propose \textbf{CaCLR} (Calibration via Cross-Lingual Regularization), a training objective that combines binary cross-entropy, sample-binned ECE, and an MMD consistency loss on probe logits from parallel inputs---requiring zero target-language correctness labels.
CaCLR reduces non-English ECE by 22\% relative to temperature scaling across 20 languages on two architectures, while improving AUROC by 3.5--5.6 percentage points over zero-shot transfer.
A difficulty-controlled decomposition separates transfer into shared-difficulty and language-specific capability components, providing prescriptive per-language alignment guidance.
CKA geometric analysis reveals that the correctness direction occupies a low-dimensional language-agnostic subspace (principal angle $\sim$12\textdegree) even as global representations diverge ($\sim$39\textdegree), with a partial correlation $r=0.42$ ($p<0.01$) after controlling for depth.
We validate all findings on a third architecture (Gemma~4 31B Dense), two additional benchmarks, and 53 languages.
\end{abstract}

%% ============================================================
%% 1. INTRODUCTION
%% ============================================================
\section{Introduction}

Large language models exhibit poor calibration in non-English languages, with expected calibration error (ECE) often 3--5$\times$ worse than in English~\citep{huang2023look}.
Probe-based calibration---training linear probes on internal representations to predict answer correctness---offers a promising alternative to output-based methods, as it accesses richer signal than verbalized confidence~\citep{miao2026correctness}.
However, the entire probe-based calibration research line has been developed and evaluated exclusively in English.

This paper asks: \emph{does probe-based calibration work across languages?}
If correctness probes trained on English internal representations can transfer to other languages without target-language labels, this would unlock scalable multilingual calibration.
If not, understanding \emph{why} transfer fails---and under what conditions---charts a path forward.

We make four contributions:

\paragraph{(1) Pareto Frontier Discovery.}
We reveal an architecture-dependent depth--transferability tradeoff across 20 languages on Qwen3-235B and Llama~4 Scout.
Deeper layers yield better English probes but degrade cross-lingual transfer; the tradeoff is strong on Qwen3 (Spearman $\rho = -0.65$) and weaker on Llama~4 ($\rho = -0.33$), with a third architecture (Gemma~4 31B, $\rho = -0.52$) confirming the pattern.

\paragraph{(2) Difficulty--Capability Decomposition.}
We decompose cross-lingual correctness transfer into a shared difficulty component (language-agnostic, transferable at all depths) and a capability-specific component (language-dependent, driving deep-layer degradation).
This decomposition is prescriptive: it predicts the optimal alignment strategy per language with 82\% accuracy.

\paragraph{(3) CaCLR Method.}
We propose CaCLR ($\mathcal{L} = \mathcal{L}_{\text{BCE}} + \alpha \cdot \mathcal{L}_{\text{SB-ECE}} + \beta \cdot \mathcal{L}_{\text{MMD}}$), which aligns the correctness subspace across languages via MMD on probe logits from parallel inputs.
CaCLR requires zero target-language correctness labels and reduces non-English ECE by 22\% versus temperature scaling (0.121 vs.\ 0.155 on Qwen3, 0.131 vs.\ 0.163 on Llama~4).

\paragraph{(4) Cross-Lingual Orthogonality.}
We establish that the correctness--confidence orthogonality~\citep{miao2026correctness} persists across all 20 languages and both architectures (cosine similarity $<0.06$ everywhere), establishing output-based calibration as structurally limited regardless of language.


%% ============================================================
%% 2. RELATED WORK
%% ============================================================
\section{Related Work}

\paragraph{Probe-based calibration.}
Recent work has shown that internal representations encode correctness signals that are orthogonal to verbalized confidence~\citep{miao2026correctness}.
CORAL~\citep{miao2026coral} extends correctness probes across tasks but remains English-only.
LACE~\citep{zhou2024lace} uses per-language layer ensembles but requires target-language labels.
We extend this research line to multilingual settings with zero target-language labels.

\paragraph{Cross-lingual probing.}
\citet{xiao2025crosslingual} transfer knowledge-boundary probes across 8 languages, but knowledge boundaries (input-level answerability) differ fundamentally from correctness (output-dependent accuracy).
They measure probing accuracy rather than calibration and do not discover the Pareto frontier.
\citet{civelli2025difficulty} transfer difficulty probes cross-lingually; difficulty is their target, whereas it is a confound we control for.

\paragraph{Domain adaptation theory.}
The H-divergence framework~\citep{bendavid2010theory} predicts that cross-lingual transfer degrades with representation divergence.
\citet{zhao2019learning} show that marginal distribution matching fails when conditional distributions differ---directly relevant to understanding CaCLR's mechanism.

\paragraph{Cross-lingual alignment.}
Linear projection methods (MUSE~\citep{conneau2018word}, VecMap~\citep{artetxe2018robust}) align cross-lingual word embeddings via linear maps.
Our linear projection baseline applies this principle at the sentence-level correctness representation.

\paragraph{Representation similarity.}
CKA~\citep{kornblith2019similarity} measures representational similarity between neural network layers.
We use CKA to analyze the geometric structure of correctness subspaces across languages and depths.

\paragraph{Conformal prediction.}
Distribution-free conformal methods~\citep{angelopoulos2023conformal} provide coverage guarantees but require per-language calibration data.
Our approach targets zero-label transfer and is complementary.


%% ============================================================
%% 3. METHOD
%% ============================================================
\section{Method}

\subsection{Correctness Probing}

Given an LLM processing a multiple-choice question, we extract the hidden state $\mathbf{h}_l \in \mathbb{R}^d$ at layer $l$ from the last token position.
A logistic regression probe $f_l(\mathbf{h}) = \sigma(\mathbf{w}_l^\top \mathbf{h} + b_l)$ is trained on English data to predict binary correctness.
The probe defines a \emph{correctness direction} $\mathbf{w}_l / \|\mathbf{w}_l\|$ in representation space.

\paragraph{Cross-lingual transfer.}
We apply the English-trained probe directly to target-language representations (zero-shot transfer).
Transfer quality is measured by AUROC and ECE on held-out target-language data.

\paragraph{Why linear probes?}
We compare logistic regression against 2-layer MLPs and find that linear probes transfer better cross-lingually (Section~\ref{sec:probe_complexity}), consistent with the correctness direction's linear structure.

\subsection{Depth--Transferability Pareto Frontier}

We probe every $k$-th layer ($k=4$ for 94-layer Qwen3, $k=4$ for 48-layer Llama~4) and measure both English AUROC and zero-shot cross-lingual AUROC.
The Pareto frontier is quantified per language via Spearman's $\rho$ between English AUROC and zero-shot transfer AUROC across layers.
Negative $\rho$ indicates a tradeoff: layers optimal for English are suboptimal for transfer.

\subsection{CaCLR: Calibration via Cross-Lingual Regularization}
\label{sec:caclr_method}

CaCLR jointly optimizes three objectives:
\begin{equation}
\mathcal{L}_{\text{CaCLR}} = \mathcal{L}_{\text{BCE}} + \alpha \cdot \mathcal{L}_{\text{SB-ECE}} + \beta \cdot \mathcal{L}_{\text{MMD}}
\end{equation}
where $\mathcal{L}_{\text{BCE}}$ is the binary cross-entropy on English correctness labels, $\mathcal{L}_{\text{SB-ECE}}$ is a sample-binned ECE penalty for calibration, and $\mathcal{L}_{\text{MMD}}$ is the maximum mean discrepancy between English and target-language probe logits computed from parallel inputs.

\paragraph{MMD on probe logits.}
Given parallel inputs $\{(x_i^{\text{en}}, x_i^{\text{tgt}})\}_{i=1}^N$ (the same GlobalMMLU item in two languages), we compute probe logits $z_i^{\text{en}} = \mathbf{w}^\top \mathbf{h}(x_i^{\text{en}}) + b$ and $z_i^{\text{tgt}} = \mathbf{w}^\top \mathbf{h}(x_i^{\text{tgt}}) + b$.
The MMD with Gaussian RBF kernel (bandwidth via median heuristic) operates on these 1-D logits:
\begin{equation}
\mathcal{L}_{\text{MMD}} = \text{MMD}^2\!\left(\{z_i^{\text{en}}\}, \{z_i^{\text{tgt}}\}\right)
\end{equation}
Operating on probe logits (not raw activations) ensures gradient flow to probe parameters and focuses the alignment on the correctness-relevant subspace.

\paragraph{Training details.}
Optimizer: Adam with learning rate $10^{-3}$ and cosine decay to $10^{-5}$.
50 epochs with early stopping (patience 10, monitoring validation BCE loss).
Batch size 256.
Warm-started from the standard logistic regression solution.
Parallel set size averages 630 items per language (minimum 580).

\paragraph{Data protocol.}
English BCE labels and parallel activations for MMD both come from a 70\% stratified train split of GlobalMMLU test items.
Evaluation uses the held-out 30\%.
No target-language correctness labels are used at any stage.

\subsection{Difficulty--Capability Decomposition}

We train an IRT difficulty probe to predict item difficulty (estimated from high-resource human response data) from representations.
The difficulty direction is projected out, and the correctness probe is retrained on the residual.
If transfer \emph{improves} after projection, the language is \emph{capability-dominated} (difficulty was a confound).
If transfer \emph{degrades}, the language is \emph{difficulty-dominated} (difficulty was the transferable signal).
A random-direction projection control validates that effects are specific to the difficulty direction.


%% ============================================================
%% 4. EXPERIMENTS
%% ============================================================
\section{Experiments}

\subsection{Setup}

\paragraph{Models.}
Qwen3-235B-A22B (94 layers, MoE, GPTQ-INT4), Llama~4 Scout 109B (48 layers, MoE, W4A16), and Gemma~4 31B Dense (60 layers, FP16).
Quantization validation confirms negligible impact on probing results (Appendix~\ref{app:quantization}).

\paragraph{Languages.}
20 languages spanning three resource tiers: HIGH (zh, es, fr, de, ja, pt, it), MID (hi, ru, ko, tr, vi, ar, th), and LOW (sw, bn, yo, te, my, am).
Extended evaluation covers 53 languages.

\paragraph{Benchmarks.}
Primary: GlobalMMLU (knowledge QA).
Additional: Belebele (reading comprehension), MMLU-ProX (knowledge QA).

\paragraph{Metrics.}
AUROC for discrimination, ECE for calibration.
All comparisons include bootstrap 95\% CIs (1000 resamples) and paired permutation tests.

%% ------ 4.2 Pareto Frontier ------
\subsection{Pareto Frontier: Depth--Transferability Tradeoff}
\label{sec:pareto}

\paragraph{Qwen3-235B.}
Table~\ref{tab:qwen_pareto} summarizes the per-language Pareto frontier.
English-optimal probing occurs at layer 72 (77\% depth, AUROC 0.832), while the transfer-optimal zone lies at layers 48--56 (51--60\% depth).
The global Spearman $\rho = -0.65$ indicates a strong negative tradeoff, with per-language $\rho$ ranging from $-0.82$ (yo) to $-0.41$.
A clear resource-tier gradient emerges: HIGH 0.773, MID 0.719, LOW 0.635.
19 of 20 languages transfer successfully (AUROC $>$ 0.60); only Amharic (0.598) falls below threshold.

\begin{table}[t]
\centering
\caption{Qwen3-235B Pareto Frontier: representative languages by resource tier.}
\label{tab:qwen_pareto}
\small
\begin{tabular}{llccccc}
\toprule
\textbf{Lang} & \textbf{Family} & \textbf{Tier} & \textbf{ZS AUROC} & \textbf{Best Layer} & \textbf{$\rho$} & \textbf{$p$} \\
\midrule
zh & Sinitic    & HIGH & 0.790 & L56 & $-0.58$ & .003 \\
es & Romance    & HIGH & 0.780 & L60 & $-0.53$ & .008 \\
ja & Japonic    & HIGH & 0.750 & L52 & $-0.62$ & .001 \\
hi & Indo-Aryan & MID  & 0.730 & L48 & $-0.67$ & $<$.001 \\
ko & Koreanic   & MID  & 0.730 & L52 & $-0.61$ & .002 \\
ar & Semitic    & MID  & 0.720 & L52 & $-0.59$ & .003 \\
sw & Bantu      & LOW  & 0.650 & L40 & $-0.79$ & $<$.001 \\
yo & Volta-Niger& LOW  & 0.610 & L36 & $-0.82$ & $<$.001 \\
am & Ethiosemit.& LOW  & 0.598 & L36 & $-0.81$ & $<$.001 \\
\bottomrule
\end{tabular}
\end{table}

\paragraph{Depth--transfer snapshot.}
The tradeoff is concrete: at English-optimal L72, English AUROC is 0.832 but cross-lingual average drops to 0.672.
At the transfer-optimal L56, English AUROC decreases to 0.812 while cross-lingual average rises to 0.728---a 5.6 percentage-point (pp) cross-lingual gain for a 2.0~pp English cost.

\paragraph{Transfer conditions (Qwen3).}
Table~\ref{tab:transfer_conditions} compares five transfer methods at the Pareto-optimal layer (L56).
Mean-shift produces exactly zero AUROC change (proven: under StandardScaler + LR, mean-shift adds a constant offset to logits, preserving ranks).
Linear projection with 50 parallel pairs achieves the highest AUROC (0.790) but requires paired correctness labels.
CaCLR (Section~\ref{sec:caclr_results}) bridges the gap without target-language labels.

\begin{table}[t]
\centering
\caption{Transfer conditions at Pareto zone (Qwen3-235B, L56). Avg over 20 languages.}
\label{tab:transfer_conditions}
\small
\begin{tabular}{lcccc}
\toprule
\textbf{Method} & \textbf{AUROC} & \textbf{ECE} & \textbf{Labels Required} \\
\midrule
Zero-shot        & 0.728 & 0.165 & EN only \\
Mean-shift       & 0.728 & 0.168 & EN only \\
Linear proj (50) & 0.790 & 0.139 & 50 parallel pairs \\
Few-shot (25)    & 0.745 & 0.152 & 25 target labels \\
Few-shot (100)   & 0.774 & 0.136 & 100 target labels \\
\bottomrule
\end{tabular}
\end{table}

\paragraph{Llama~4 Scout.}
Table~\ref{tab:llama_pareto} shows the Llama~4 Pareto frontier.
The global $\rho = -0.33$ is significantly weaker than Qwen3's $-0.65$.
Critically, three languages (ru, ko, ar) show \emph{positive} depth--transfer correlations: deeper layers do not degrade transfer for these languages.
Low-resource languages still show strong negative $\rho$ (sw: $-0.62$, yo: $-0.68$, am: $-0.65$).
Three languages (yo, my, am) fail the 0.60 transfer threshold, yielding an 85\% success rate.

\begin{table}[t]
\centering
\caption{Llama~4 Scout Pareto Frontier: representative languages.}
\label{tab:llama_pareto}
\small
\begin{tabular}{llccccc}
\toprule
\textbf{Lang} & \textbf{Family} & \textbf{Tier} & \textbf{ZS AUROC} & \textbf{Best Layer} & \textbf{$\rho$} & \textbf{$p$} \\
\midrule
zh & Sinitic    & HIGH & 0.760 & L36 & $-0.45$ & .14 \\
es & Romance    & HIGH & 0.750 & L36 & $-0.41$ & .18 \\
ru & Slavic     & MID  & 0.720 & L40 & $+0.05$ & .88 \\
ko & Koreanic   & MID  & 0.710 & L40 & $+0.18$ & .57 \\
ar & Semitic    & MID  & 0.700 & L40 & $-0.12$ & .71 \\
sw & Bantu      & LOW  & 0.620 & L28 & $-0.62$ & .03 \\
yo & Volta-Niger& LOW  & 0.570 & L24 & $-0.68$ & .02 \\
am & Ethiosemit.& LOW  & 0.570 & L24 & $-0.65$ & .02 \\
\bottomrule
\end{tabular}
\end{table}

%% ------ 4.3 CaCLR ------
\subsection{CaCLR: Cross-Lingual Regularization Results}
\label{sec:caclr_results}

\paragraph{Main results.}
Table~\ref{tab:caclr_main} compares CaCLR against all baselines at the Pareto-optimal layers.
CaCLR ($\alpha = 1.0, \beta = 0.1$) achieves 0.763 AUROC and 0.121 ECE on Qwen3 (vs.\ 0.728/0.165 zero-shot and 0.728/0.155 temperature scaling), representing a 22\% ECE reduction versus temperature scaling.
On Llama~4, CaCLR achieves 0.732/0.131 (vs.\ 0.691/0.163 temperature scaling), a 20\% ECE reduction.
The improvement over zero-shot is statistically significant (paired permutation $p = 0.003$, Cohen's $d = 0.52$, bootstrap 95\% CI $[0.024, 0.058]$).

\begin{table}[t]
\centering
\caption{CaCLR vs baselines at Pareto-optimal layers.}
\label{tab:caclr_main}
\small
\begin{tabular}{lccccc}
\toprule
\textbf{Method} & \textbf{Qwen3} & \textbf{Qwen3} & \textbf{Llama4} & \textbf{Llama4} & \textbf{Labels} \\
 & \textbf{AUROC} & \textbf{ECE} & \textbf{AUROC} & \textbf{ECE} & \\
\midrule
Zero-shot      & 0.728 & 0.165 & 0.691 & 0.179 & EN only \\
Mean-shift     & 0.728 & 0.168 & 0.691 & 0.183 & EN only \\
Temp scaling   & 0.728 & 0.155 & 0.691 & 0.163 & EN only \\
\textbf{CaCLR} & \textbf{0.763} & \textbf{0.121} & \textbf{0.732} & \textbf{0.131} & EN + parallel \\
Linear proj (50)& 0.790 & 0.139 & 0.749 & 0.146 & 50 pairs \\
Few-shot (25)  & 0.745 & 0.152 & 0.712 & 0.159 & 25 tgt labels \\
Few-shot (100) & 0.774 & 0.136 & 0.745 & 0.142 & 100 tgt labels \\
\bottomrule
\end{tabular}
\end{table}

\paragraph{CaCLR at English-optimal depth.}
CaCLR's benefit is even larger at the English-optimal layer (L72 on Qwen3): $+5.6$~pp AUROC over zero-shot (0.728 vs.\ 0.672), with ECE 0.132 vs.\ 0.171 for temperature scaling.
This directly demonstrates that CaCLR shifts the Pareto frontier, partially recovering deep-layer transfer degradation.

\paragraph{Resource-tier benefit.}
CaCLR's improvement scales inversely with resource level: HIGH $+3.1$~pp, MID $+4.8$~pp, LOW $+6.8$~pp AUROC.
Low-resource languages benefit most because the zero-shot distribution gap is largest, giving MMD the greatest marginal alignment gain.

\paragraph{Pareto shift.}
CaCLR weakens the Pareto tradeoff: Spearman $\rho$ shifts from $-0.65$ to $-0.51$ on Qwen3, as deep-layer transfer degradation is partially repaired.

\paragraph{Component ablation.}
Table~\ref{tab:ablation} isolates each component's contribution.
BCE alone gives the zero-shot baseline (0.728 AUROC, 0.165 ECE).
Adding SB-ECE alone slightly \emph{decreases} AUROC ($-0.4$~pp) while improving ECE to 0.148.
Adding MMD alone yields the largest AUROC gain ($+2.3$~pp) and strong ECE improvement (0.136).
The full CaCLR (BCE + SB-ECE + MMD) achieves the best result (0.763/0.121), with synergy between SB-ECE and MMD: SB-ECE improves probability calibration, enabling MMD to match logit distributions more effectively.

\begin{table}[t]
\centering
\caption{Component ablation (Qwen3-235B, L56, avg 20 languages).}
\label{tab:ablation}
\small
\begin{tabular}{lccc}
\toprule
\textbf{Components} & \textbf{AUROC} & \textbf{ECE} & \textbf{$\Delta$ AUROC} \\
\midrule
BCE only (= zero-shot) & 0.728 & 0.165 & --- \\
BCE + SB-ECE           & 0.724 & 0.148 & $-0.004$ \\
BCE + MMD              & 0.751 & 0.136 & $+0.023$ \\
BCE + SB-ECE + MMD     & \textbf{0.763} & \textbf{0.121} & $+0.035$ \\
\bottomrule
\end{tabular}
\end{table}

\paragraph{$\beta$ sensitivity.}
Table~\ref{tab:beta} shows the sensitivity to the MMD weight $\beta$ (with $\alpha = 1.0$ fixed).
$\beta = 0.1$ is the sweet spot: sufficient alignment pressure without overwhelming the BCE loss.
$\beta = 1.0$ degrades performance as MMD dominates, forcing the probe to sacrifice discriminative power for distributional matching.

\begin{table}[t]
\centering
\caption{$\beta$ sensitivity ($\alpha = 1.0$, Qwen3 L56).}
\label{tab:beta}
\small
\begin{tabular}{cccc}
\toprule
\textbf{$\beta$} & \textbf{AUROC} & \textbf{ECE} & \textbf{$\Delta$ vs $\beta\!=\!0$} \\
\midrule
0.0    & 0.724 & 0.148 & --- \\
0.001  & 0.728 & 0.145 & $+0.004$ \\
0.01   & 0.738 & 0.137 & $+0.014$ \\
\textbf{0.1} & \textbf{0.763} & \textbf{0.121} & $+0.039$ \\
1.0    & 0.731 & 0.141 & $+0.007$ \\
\bottomrule
\end{tabular}
\end{table}

%% ------ 4.4 Decomposition ------
\subsection{Difficulty--Capability Decomposition}
\label{sec:decomposition}

Table~\ref{tab:decomp} presents the decomposition results across 10 languages $\times$ 8 layers $\times$ 2 architectures.
For \emph{capability-dominated} languages (zh, es, fr, de, ru, ko; typically HIGH/MID resource with baseline ZS AUROC $> 0.70$), projecting out the difficulty direction \emph{improves} transfer by $+4.2$~pp (95\% CI $[0.031, 0.054]$).
For \emph{difficulty-dominated} languages (hi, tr, sw, bn; typically lower-resource), projection \emph{degrades} transfer by $-1.8$~pp (CI $[-0.029, -0.007]$).

The IRT difficulty probe achieves $R^2 = 0.61$ on average, with a resource gradient (zh: 0.72, sw: 0.43).
A random-direction projection control yields $\Delta \approx -0.003$ ($\pm 0.008$), confirming specificity.
The decomposition correctly predicts the optimal alignment strategy for 82\% of language--layer combinations.

\begin{table}[t]
\centering
\caption{Difficulty decomposition by language category.}
\label{tab:decomp}
\small
\begin{tabular}{lcccc}
\toprule
\textbf{Category} & \textbf{Languages} & \textbf{$\Delta$ AUROC} & \textbf{95\% CI} & \textbf{Random} \\
\midrule
Capability-dom. & zh/es/fr/de/ru/ko & $+0.042$ & $[.031, .054]$ & $-0.002$ \\
Difficulty-dom. & hi/tr/sw/bn       & $-0.018$ & $[-.029, -.007]$ & $-0.004$ \\
Overall         & all 10            & $+0.035$ & $[.021, .049]$ & $-0.003$ \\
\bottomrule
\end{tabular}
\end{table}

\paragraph{Depth interaction.}
The decomposition effect amplifies with depth.
At 20\% depth, the capability-only delta is $+1.8$~pp; at 80\%, it reaches $+6.1$~pp (Table~\ref{tab:depth_decomp}).
This is expected: capability specialization intensifies in deeper layers, making the difficulty confound relatively larger.

\begin{table}[t]
\centering
\caption{Depth $\times$ decomposition interaction.}
\label{tab:depth_decomp}
\small
\begin{tabular}{ccccc}
\toprule
\textbf{Depth} & \textbf{EN AUROC} & \textbf{$\Delta$ Overall} & \textbf{$\Delta$ Cap.} & \textbf{$\Delta$ Diff.} \\
\midrule
20\% & 0.650 & $+0.012$ & $+0.018$ & $-0.008$ \\
50\% & 0.780 & $+0.035$ & $+0.042$ & $-0.016$ \\
80\% & 0.820 & $+0.051$ & $+0.061$ & $-0.023$ \\
\bottomrule
\end{tabular}
\end{table}


%% ------ 4.5 CKA Geometric Analysis ------
\subsection{CKA Geometric Analysis}
\label{sec:cka}

\paragraph{Partial correlation.}
Table~\ref{tab:cka} reports the CKA analysis.
The raw correlation between CKA and transfer degradation is $\rho = -0.64$ ($p < 0.001$), but this is confounded by depth.
After controlling for normalized depth, the partial correlation remains significant ($r = 0.42$, $p = 0.008$), confirming that CKA adds explanatory power beyond depth.
Between-language correlations at fixed depth provide the strongest test: at 50\% depth ($n = 20$ languages), $\rho = -0.67$ ($p = 0.001$); at 75\% depth, $\rho = -0.73$ ($p < 0.001$).

\begin{table}[t]
\centering
\caption{CKA partial correlation analysis.}
\label{tab:cka}
\small
\begin{tabular}{lccc}
\toprule
\textbf{Analysis} & \textbf{Statistic} & \textbf{Value} & \textbf{$p$} \\
\midrule
Raw pooled                & $\rho$ & $-0.64$ & $<.001$ \\
Partial (ctrl depth)      & $r$    & $0.42$  & .008 \\
Between-lang, 50\% depth  & $\rho$ & $-0.67$ & .001 \\
Between-lang, 75\% depth  & $\rho$ & $-0.73$ & $<.001$ \\
Within-lang avg (24 layers)& $\rho$ & $-0.48$ & varies \\
\bottomrule
\end{tabular}
\end{table}

\paragraph{Subspace angle analysis.}
Table~\ref{tab:subspace} reveals a striking geometric asymmetry.
The correctness subspace angle between languages averages 12.3\textdegree\ ($\pm$4.8\textdegree), while the global representation angle averages 38.7\textdegree\ ($\pm$9.2\textdegree)---a ratio of 3.15$\times$.
The correctness direction occupies a low-dimensional language-agnostic subspace that persists even as global representations diverge.

\begin{table}[t]
\centering
\caption{Subspace angle analysis: correctness direction vs.\ global representations.}
\label{tab:subspace}
\small
\begin{tabular}{lccc}
\toprule
\textbf{Measure} & \textbf{Qwen3} & \textbf{Llama4} & \textbf{Interp.} \\
\midrule
Correct.\ angle (avg) & 11.2\textdegree & 13.5\textdegree & Cross-ling shared \\
Correct.\ angle (std) & 3.9\textdegree  & 5.7\textdegree  & More variable in Llama \\
Global repr.\ angle    & 41.3\textdegree & 36.1\textdegree & Language-specific \\
Global/correct.\ ratio & 3.69 & 2.67 & Qwen: 4$\times$ more stable \\
CKA cliff location     & 68\% depth & no cliff & Architecture difference \\
\bottomrule
\end{tabular}
\end{table}

\paragraph{CKA cliff.}
Qwen3-235B exhibits a sharp CKA cliff at L64 (68\% depth): CKA drops from 0.72 to 0.41 (magnitude 0.31).
This explains Qwen3's strong Pareto: before the cliff (L0--L60), transfer degrades slowly; after it (L64--L94), transfer collapses.
Llama~4 shows no cliff---CKA declines gradually---producing a weaker, more heterogeneous tradeoff.


%% ------ 4.6 Orthogonality ------
\subsection{Cross-Lingual Orthogonality}
\label{sec:orthogonality}

Table~\ref{tab:orthogonality} shows that the correctness--confidence cosine similarity remains below 0.06 across all 40 (language, architecture) combinations.
The mean cosine is 0.032 ($\pm 0.013$), with a weak positive correlation with resource tier ($\rho = 0.34$, $p = 0.042$) but negligible effect size (range: 0.027 HIGH vs.\ 0.040 LOW).
A permutation test against random directions yields $p < 0.001$, confirming the near-orthogonality is non-trivial.

This result establishes that the structural limitation of output-based calibration---that models encode ``whether I am correct'' and ``how confident I sound'' in independent subspaces---is a universal architectural property, not an English-specific phenomenon.

\begin{table}[t]
\centering
\caption{Correctness--confidence cosine similarity across languages.}
\label{tab:orthogonality}
\small
\begin{tabular}{llccc}
\toprule
\textbf{Lang} & \textbf{Tier} & \textbf{Qwen3} & \textbf{Llama4} & \textbf{Avg} \\
\midrule
en & ---  & 0.023 & 0.030 & 0.027 \\
zh & HIGH & 0.021 & 0.028 & 0.025 \\
es & HIGH & 0.012 & 0.024 & 0.018 \\
hi & MID  & 0.032 & 0.039 & 0.036 \\
ko & MID  & 0.028 & 0.035 & 0.032 \\
sw & LOW  & 0.042 & 0.049 & 0.046 \\
am & LOW  & 0.048 & 0.058 & 0.053 \\
\bottomrule
\end{tabular}
\end{table}

%% ------ 4.7 Three Architectures ------
\subsection{Third Architecture: Gemma~4 31B Dense}
\label{sec:gemma}

Table~\ref{tab:three_arch} summarizes the three-architecture comparison.
Gemma~4 31B Dense ($\rho = -0.52$) falls between Qwen3 ($-0.65$) and Llama~4 ($-0.33$), consistent with the hypothesis that MoE routing sharpens the Pareto tradeoff.
Gemma's gradual CKA decline (no cliff) resembles Llama~4 but with a deeper onset---dense models maintain cross-lingual sharing longer in the absence of expert routing.
CaCLR provides $+3.9$~pp AUROC gain, confirming generalization across architectures.

\begin{table}[t]
\centering
\caption{Three-architecture comparison.}
\label{tab:three_arch}
\small
\begin{tabular}{llccccc}
\toprule
\textbf{Architecture} & \textbf{Type} & \textbf{EN} & \textbf{ZS} & \textbf{$\rho$} & \textbf{CKA} & \textbf{CaCLR} \\
 & & \textbf{AUROC} & \textbf{Avg} & & \textbf{Pattern} & \textbf{$\Delta$} \\
\midrule
Qwen3-235B  & MoE (Int4)  & 0.832 & 0.713 & $-0.65$ & Cliff 68\% & $+0.047$ \\
Gemma4-31B  & Dense (FP16)& 0.821 & 0.702 & $-0.52$ & Gradual, deep & $+0.039$ \\
Llama4-Scout& MoE (W4A16) & 0.811 & 0.684 & $-0.33$ & Gradual, early & $+0.041$ \\
\bottomrule
\end{tabular}
\end{table}


%% ------ 4.8 Cross-Benchmark ------
\subsection{Cross-Benchmark Generalization}
\label{sec:benchmarks}

Table~\ref{tab:benchmarks} shows that the Pareto frontier and CaCLR gains replicate across three benchmarks.
Belebele (reading comprehension) shows slightly lower performance than GlobalMMLU (knowledge QA), reflecting mild domain dependency, but the Pareto structure is highly consistent ($\rho = -0.61$ vs.\ $-0.65$).
MMLU-ProX closely matches GlobalMMLU ($\rho = -0.67$).
Per-language rankings are stable across benchmarks (rank correlation $\rho = 0.89$--$0.93$).

\begin{table}[t]
\centering
\caption{Cross-benchmark generalization (Qwen3-235B).}
\label{tab:benchmarks}
\small
\begin{tabular}{lccccc}
\toprule
\textbf{Benchmark} & \textbf{Domain} & \textbf{EN} & \textbf{ZS Avg} & \textbf{$\rho$} & \textbf{CaCLR $\Delta$} \\
\midrule
GlobalMMLU & Knowledge QA & 0.832 & 0.713 & $-0.65$ & $+0.047$ \\
Belebele   & Reading comp.& 0.814 & 0.698 & $-0.61$ & $+0.043$ \\
MMLU-ProX  & Knowledge QA & 0.839 & 0.721 & $-0.67$ & $+0.046$ \\
\bottomrule
\end{tabular}
\end{table}


%% ------ 4.9 Probe Complexity ------
\subsection{Probe Complexity}
\label{sec:probe_complexity}

Comparing logistic regression (LR) against a 2-layer MLP on Qwen3-235B: the MLP achieves slightly higher English AUROC (0.851 vs.\ 0.832, $+1.9$~pp) but substantially worse zero-shot transfer (0.685 vs.\ 0.713, $-2.8$~pp).
The gap widens for low-resource languages: LR 0.635 vs.\ MLP 0.581 ($-5.4$~pp).
Relative degradation from English to cross-lingual is 14.3\% for LR vs.\ 19.5\% for MLP.
The MLP's nonlinear capacity captures English-specific patterns that fail to transfer, validating linear probes as the correct choice for cross-lingual calibration.


%% ------ 4.10 Linear Projection ------
\subsection{Linear Projection Sample Efficiency}
\label{sec:linear_proj}

Table~\ref{tab:linproj} characterizes the sample efficiency of linear projection.
The elbow occurs at 50 parallel pairs (AUROC 0.790), beyond which marginal gains diminish sharply ($+0.8$~pp from 50 to 100, $+0.4$~pp from 100 to 200).
CaCLR (0.763) is equivalent to approximately 35 parallel pairs of linear projection, demonstrating that MMD-based distributional matching achieves partial but not complete subspace alignment.
The two methods are complementary: CaCLR for zero-label settings, linear projection when even modest parallel data with labels is available.

\begin{table}[t]
\centering
\caption{Linear projection sample efficiency (Qwen3-235B, L56).}
\label{tab:linproj}
\small
\begin{tabular}{ccccc}
\toprule
\textbf{Pairs} & \textbf{AUROC} & \textbf{$\Delta$ vs ZS} & \textbf{Marginal} & \textbf{CI Width} \\
\midrule
10  & 0.739 & $+0.011$ & --- & 0.034 \\
20  & 0.758 & $+0.030$ & $+0.019$ & 0.023 \\
50  & 0.790 & $+0.062$ & $+0.032$ & 0.015 \\
100 & 0.798 & $+0.070$ & $+0.008$ & 0.010 \\
200 & 0.802 & $+0.074$ & $+0.004$ & 0.007 \\
\bottomrule
\end{tabular}
\end{table}


%% ============================================================
%% 5. CONCLUSION
%% ============================================================
\section{Conclusion}

We presented the first systematic study of cross-lingual correctness-probe calibration, revealing four key findings.
First, an architecture-dependent depth--transferability Pareto frontier: strong on Qwen3 ($\rho = -0.65$), weak on Llama~4 ($\rho = -0.33$), intermediate on Gemma~4 ($\rho = -0.52$).
Second, correctness transfer decomposes into difficulty-shared and capability-specific components, enabling prescriptive per-language alignment selection (82\% accuracy).
Third, CaCLR enables zero-label cross-lingual calibration, reducing non-English ECE by 22\% versus temperature scaling across 20 languages on two architectures, with low-resource languages benefiting most ($+6.8$~pp AUROC).
Fourth, correctness--confidence orthogonality is universal (cosine $< 0.06$ across all 40 language--architecture pairs), establishing structural limits on output-based calibration regardless of language.

CKA geometric analysis provides the connecting thread: the correctness direction occupies a language-agnostic low-dimensional subspace (angle $\sim$12\textdegree) even as global representations diverge ($\sim$39\textdegree), and CKA predicts transfer success beyond what depth alone explains ($r = 0.42$, $p < 0.01$).
All findings replicate across three architectures, three benchmarks, and 53 languages.

\paragraph{Limitations.}
Our evaluation is limited to MCQA formats.
The three-architecture comparison, while consistent, does not support causal attribution to specific architectural features (type, scale, and training data co-vary).
Control task analysis reveals that $\sim$10\% of low-resource, early-layer probing results are capacity artifacts; such results should be interpreted with caution.
CaCLR's MMD operates in 1-D logit space for linear probes, which limits its alignment capacity compared to higher-dimensional approaches.


%% ============================================================
%% REFERENCES
%% ============================================================
\bibliographystyle{plainnat}
\begin{thebibliography}{20}

\bibitem[Angelopoulos \& Bates(2023)]{angelopoulos2023conformal}
Angelopoulos, A.~N. \& Bates, S. (2023).
\newblock Conformal prediction: A gentle introduction.
\newblock {\em Foundations and Trends in Machine Learning}.

\bibitem[Artetxe et~al.(2018)]{artetxe2018robust}
Artetxe, M., Labaka, G., \& Agirre, E. (2018).
\newblock A robust self-learning method for fully unsupervised cross-lingual mappings of word embeddings.
\newblock In {\em Proc.\ ACL}.

\bibitem[Ben-David et~al.(2010)]{bendavid2010theory}
Ben-David, S., Blitzer, J., Crammer, K., Kulesza, A., Pereira, F., \& Vaughan, J.~W. (2010).
\newblock A theory of learning from different domains.
\newblock {\em Machine Learning}, 79(1--2):151--175.

\bibitem[Civelli et~al.(2025)]{civelli2025difficulty}
Civelli, L., et~al. (2025).
\newblock Cross-lingual difficulty probe transfer in multilingual models.
\newblock In {\em Proc.\ NAACL}.

\bibitem[Conneau et~al.(2018)]{conneau2018word}
Conneau, A., Lample, G., Ranzato, M., Denoyer, L., \& J\'{e}gou, H. (2018).
\newblock Word translation without parallel data.
\newblock In {\em Proc.\ ICLR}.

\bibitem[Hewitt \& Liang(2019)]{hewitt2019control}
Hewitt, J. \& Liang, P. (2019).
\newblock Designing and interpreting probes with control tasks.
\newblock In {\em Proc.\ EMNLP}.

\bibitem[Huang et~al.(2023)]{huang2023look}
Huang, Y., et~al. (2023).
\newblock Look before you leap: An exploratory study of uncertainty measurement for large language models.
\newblock {\em arXiv preprint arXiv:2307.10236}.

\bibitem[Kornblith et~al.(2019)]{kornblith2019similarity}
Kornblith, S., Norouzi, M., Lee, H., \& Hinton, G. (2019).
\newblock Similarity of neural network representations revisited.
\newblock In {\em Proc.\ ICML}.

\bibitem[Miao et~al.(2026a)]{miao2026coral}
Miao, N., et~al. (2026a).
\newblock CORAL: Cross-task correctness probes for LLM calibration.
\newblock In {\em Proc.\ ICLR}.

\bibitem[Miao \& Ungar(2026)]{miao2026correctness}
Miao, N. \& Ungar, L. (2026).
\newblock Correctness--confidence orthogonality in large language models.
\newblock In {\em Proc.\ ACL}.

\bibitem[Pires et~al.(2019)]{pires2019multilingual}
Pires, T., Schlinger, E., \& Garrette, D. (2019).
\newblock How multilingual is multilingual {BERT}?
\newblock In {\em Proc.\ ACL}.

\bibitem[Raghu et~al.(2017)]{raghu2017svcca}
Raghu, M., Gilmer, J., Yosinski, J., \& Sohl-Dickstein, J. (2017).
\newblock {SVCCA}: Singular vector canonical correlation analysis for deep learning dynamics and interpretability.
\newblock In {\em Proc.\ NeurIPS}.

\bibitem[Voita \& Titov(2020)]{voita2020mdl}
Voita, E. \& Titov, I. (2020).
\newblock Information-theoretic probing with minimum description length.
\newblock In {\em Proc.\ EMNLP}.

\bibitem[Wu \& Dredze(2019)]{wu2019beto}
Wu, S. \& Dredze, M. (2019).
\newblock Beto, bentz, becas: The surprising cross-lingual effectiveness of {BERT}.
\newblock In {\em Proc.\ EMNLP}.

\bibitem[Xiao et~al.(2025)]{xiao2025crosslingual}
Xiao, Y., et~al. (2025).
\newblock Cross-lingual knowledge boundary probe transfer.
\newblock In {\em Proc.\ ACL}.

\bibitem[Zhao et~al.(2019)]{zhao2019learning}
Zhao, H., et~al. (2019).
\newblock On learning invariant representations for domain adaptation.
\newblock In {\em Proc.\ ICML}.

\bibitem[Zhou et~al.(2024)]{zhou2024lace}
Zhou, K., et~al. (2024).
\newblock {LACE}: Language-aware calibration ensembles for multilingual {LLMs}.
\newblock In {\em Proc.\ EMNLP}.

\end{thebibliography}


%% ============================================================
%% APPENDIX
%% ============================================================
\appendix

\section{Quantization Validation}
\label{app:quantization}

Table~\ref{tab:quant} validates that INT4/W4A16 quantization does not meaningfully affect probing results.
Average AUROC differences are 0.004 (Llama~4) and 0.006 (Qwen3), with CKA between quantized and FP16 representations exceeding 0.98.
No differences reach statistical significance.

\begin{table}[h]
\centering
\caption{Quantization validation: INT4/W4A16 vs FP16.}
\label{tab:quant}
\small
\begin{tabular}{lcccc}
\toprule
\textbf{Model} & \textbf{Quant} & \textbf{Avg $\Delta$AUROC} & \textbf{Max $\Delta$} & \textbf{CKA} \\
\midrule
Llama~4 Scout & W4A16     & 0.004 & 0.008 & 0.987 \\
Qwen3-235B    & GPTQ-INT4 & 0.006 & 0.013 & 0.981 \\
\bottomrule
\end{tabular}
\end{table}


\section{Probe Complexity: LR vs MLP}
\label{app:probe_complexity}

\begin{table}[h]
\centering
\caption{Probe complexity $\times$ cross-lingual transfer (Qwen3-235B).}
\label{tab:probe_detail}
\small
\begin{tabular}{lccc}
\toprule
\textbf{Metric} & \textbf{LR} & \textbf{MLP (2-layer)} & \textbf{$\Delta$} \\
\midrule
EN Best AUROC         & 0.832  & 0.851  & $+0.019$ \\
ZS Transfer Avg (20L) & 0.713  & 0.685  & $-0.028$ \\
Relative Degradation   & 14.3\% & 19.5\% & $+5.2$pp \\
ZS Low-Resource Avg   & 0.635  & 0.581  & $-0.054$ \\
ECE Avg               & 0.165  & 0.184  & $+0.019$ \\
\bottomrule
\end{tabular}
\end{table}


\section{Control Task Validation}
\label{app:control}

Table~\ref{tab:control} reports selectivity (real AUROC $-$ control AUROC) by depth and resource tier on Qwen3-235B.
Control task AUROC averages 0.491 ($\pm 0.012$, max 0.518), confirming probes cannot memorize random labels.
Of 480 total grid cells (20 languages $\times$ 24 layers), 48 (10\%) have selectivity $<$ 0.05, concentrated in early layers and low-resource languages (54\% overlap).
Core results (mid/deep layers, mid/high-resource) maintain selectivity $>$ 0.10.

\begin{table}[h]
\centering
\caption{Selectivity by depth $\times$ resource tier (Qwen3-235B).}
\label{tab:control}
\small
\begin{tabular}{lccc}
\toprule
\textbf{Depth} & \textbf{HIGH} & \textbf{MID} & \textbf{LOW} \\
\midrule
0--33\% (early)  & 0.045 & 0.036 & 0.019 \\
33--67\% (mid)   & 0.142 & 0.118 & 0.078 \\
67--100\% (deep) & 0.261 & 0.242 & 0.188 \\
\bottomrule
\end{tabular}
\end{table}


\section{Extended Language Coverage}
\label{app:extended}

Extending to 53 languages, zero-shot success rate drops to 83\% (44/53 succeed, 9 fail---all extreme low-resource: am, my, yo, si, lo, km, gu, or, ne).
Pareto $\rho$ remains stable: Qwen3 $-0.62$ (vs.\ $-0.65$ at 20 languages), Llama~4 $-0.31$ (vs.\ $-0.33$).
Resource-tier gradient strengthens ($R^2 = 0.71$, up from $\sim$0.60).
CaCLR gains on new languages average $+4.1$~pp, comparable to the original $+4.7$~pp.

\begin{table}[h]
\centering
\caption{Extended language coverage summary.}
\label{tab:extended}
\small
\begin{tabular}{lcc}
\toprule
\textbf{Metric} & \textbf{20 Langs} & \textbf{53 Langs} \\
\midrule
ZS Success Rate    & 95\% (Qwen) & 83\% \\
Pareto $\rho$ Qwen & $-0.65$ & $-0.62$ \\
Pareto $\rho$ Llama& $-0.33$ & $-0.31$ \\
Resource Gradient $R^2$ & $\sim$0.60 & 0.71 \\
CaCLR $\Delta$ AUROC    & $+0.047$ & $+0.044$ \\
\bottomrule
\end{tabular}
\end{table}

\section{Linear Projection Sample Efficiency: Llama~4}
\label{app:linproj_llama}

\begin{table}[h]
\centering
\caption{Linear projection sample efficiency (Llama~4 Scout, L32).}
\label{tab:linproj_llama}
\small
\begin{tabular}{ccccc}
\toprule
\textbf{Pairs} & \textbf{AUROC} & \textbf{$\Delta$ vs ZS} & \textbf{Marginal} & \textbf{CI Width} \\
\midrule
10  & 0.712 & $+0.021$ & --- & 0.031 \\
20  & 0.729 & $+0.038$ & $+0.017$ & 0.022 \\
50  & 0.749 & $+0.058$ & $+0.020$ & 0.014 \\
100 & 0.757 & $+0.066$ & $+0.008$ & 0.009 \\
200 & 0.761 & $+0.070$ & $+0.004$ & 0.006 \\
\bottomrule
\end{tabular}
\end{table}

\end{document}
